\section{从广告入口到智力劳动：Agent 的商业模式重估}

商业模式讨论从一个很经典的对照开始：互联网时代的大模式主要是广告、订阅和交易抽佣。广告卖的是注意力，订阅卖的是工具效率，抽佣卖的是交易撮合。Agent 作为入口时，会直接挑战广告，因为广告系统希望用户停留、浏览、点击，而一个好 Agent 恰恰应该减少用户停留时间，替用户完成筛选、比较和执行。

尚晏仪的判断很有启发性：互联网完成的是信息连接，广告费对应的是信息分发；Agent 输出的是智力劳动，而人类社会里智力劳动通常通过工资、服务费、中介费、佣金等形式付费。这就是为什么她更看好订阅和抽佣，而不是把 Agent 强行塞回广告逻辑。Agent 如果真的像秘书、顾问或代理，它的商业价值就应该来自节省时间、降低决策成本、撮合交易和承担行动结果，而不是让用户多看几秒内容。

旅行场景之所以在她的实践里成立，也正因为它天然存在 travel agent 这个历史角色。用户理解“代理”意味着什么，也能接受为规划、筛选、交易和服务付费。旅行还拥有明确交易闭环：酒店、机票、路线、签证、行程服务都能连接到抽佣或服务费。这和泛化的“超级入口”不同，后者看起来空间大，但 reward、履约和商业闭环都更模糊。

\begin{importantbox}{Agent 更像劳动力，而不是广告位}
如果一个系统的价值是替用户思考、筛选、行动和承担流程成本，那么它更接近智力劳动。广告逻辑要求占用用户注意力，Agent 逻辑要求释放用户注意力。两者并非不能结合，但默认商业重心不同。
\end{importantbox}

豆包手机的讨论，把这个问题推到更高一层。手机级 Agent 的优势非常清楚：手机掌握用户住址、偏好、账号、聊天、日程和生活习惯，因此最有机会成为生活秘书。但它同时会触碰现有 App 分发、广告推荐和交易入口的利益结构。换句话说，手机 Agent 技术上很合理，商业上却会改变很多人的饭碗。

这里有一个细节很值得抓住。尚晏仪说，如果用户直接告诉 Agent “我想去海岛”，这比短视频广告系统观察到用户在海岛视频上多停留几秒更强。因为前者是明确意图，后者只是行为推断。但明确意图不自动等于商业成功。Agent 还要把意图结构化，找到供应，比较选项，建立信任，并完成交易。广告系统擅长把模糊行为变成推荐，Agent 系统则必须把明确语言变成可靠行动。

\begin{knowledgebox}{从注意力经济到代理经济}
\begin{tabularx}{\textwidth}{>{\raggedright\arraybackslash}p{0.2\textwidth} Y Y}
\toprule
模式 & 价值来源 & 与 Agent 的关系 \\
\midrule
广告 & 用户停留、浏览、点击、曝光。 & 好 Agent 会减少停留，因此天然冲突。 \\
订阅 & 用户为效率、工具能力和可靠性付费。 & 适合高频、明确、可持续提效场景。 \\
抽佣 & 平台撮合交易并参与履约。 & 适合旅行、采购、电商、招聘等有交易闭环的场景。 \\
代理服务 & 用户为智力劳动、筛选、执行和责任边界付费。 & 最接近强 Agent 的长期形态。 \\
\bottomrule
\end{tabularx}
\end{knowledgebox}

\subsection{本章小结}

Agent 产品不能简单复刻互联网广告模式。它的价值更像智力劳动、代理服务和交易撮合。旅行之所以适合作为案例，是因为它既有传统代理角色，也有清晰交易闭环。手机 Agent 的潜力来自个人上下文，但真正难的是把上下文转成可信行动，同时处理与现有生态的商业冲突。

